%% ma: L12-B08-0193
%% lop: 12
%% chuong: 2
%% bai: 8
%% dang: 12.8.37
%% loai: TN4
%% muc: TH
%% dap_an: "C"
%% nguon: "CĐ2.6.pdf giáo viên gửi 27/10/2026 (M07, Lê Minh Khôi – Nguyễn Anh Khôi, Chương II Chủ đề 6 Ứng dụng vào các bài toán thực tế), Câu 23"
%% kien_thuc: "Bài 8 (tọa độ vectơ, độ dài), Bài 6 (tổng các lực cân bằng)"
%% trang_thai: chinh-thuc
%% ghi_chu: "Đề gốc dùng chung giả thiết cho Câu 22–23; đã chép vào từng câu và nói rõ mỗi lực có phương dọc theo chân máy"
\begin{ex}
Một chiếc máy quay được đặt trên một giá đỡ ba chân. Điểm đặt máy là $E(0;0;6)$. Các điểm tiếp xúc với mặt đất của ba chân lần lượt là
\[A_1(0;1;0),\quad A_2\left(-\dfrac{\sqrt3}2;-\dfrac12;0\right),\quad A_3\left(\dfrac{\sqrt3}2;-\dfrac12;0\right)\]
(đơn vị: mét). Trọng lực tác dụng lên máy quay là $\vec P$ có độ lớn $300$ N hướng thẳng đứng xuống dưới mặt đất $(Oxy)$. Gọi $\vec F_1$, $\vec F_2$, $\vec F_3$ là lực do ba chân máy tác dụng lên điểm $E$ để giữ hệ thống cân bằng (mỗi lực có phương dọc theo chân máy tương ứng).
Mệnh đề nào sau đây sai?
\choice{Tổng các lực tác dụng lên điểm đặt máy thỏa mãn $\vec F_1+\vec F_2+\vec F_3+\vec P=\vec 0$}{Độ lớn của lực $\vec F_1$ tác dụng lên một chân máy là $\dfrac{50\sqrt{37}}3$ N}{\True Trọng lực tác dụng lên máy quay được biểu diễn bởi vectơ $\vec P=(0;0;300)$}{Tổng hình chiếu của ba lực $\vec F_1$, $\vec F_2$, $\vec F_3$ lên mặt phẳng $(Oxy)$ bằng $\vec 0$}
\loigiai{$\vec P$ hướng xuống nên $\vec P=(0;0;-300)$: C sai.

A đúng (cân bằng). B: $\vec F_1=\dfrac{50}3(0;-1;6)$, $|\vec F_1|=\dfrac{50\sqrt{37}}3$. D: tổng thành phần nằm ngang của ba lực bằng $\vec 0$.}
\end{ex}
